% Generated draft: requires mathematical and source review.
\documentclass[11pt]{article}
\usepackage[margin=25mm]{geometry}
\usepackage{fontspec}
\setmainfont{Latin Modern Roman}
\usepackage{amsmath,amssymb,mathtools,unicode-math}
\setmathfont{Latin Modern Math}
\usepackage{enumitem,needspace,xurl,hyperref}
\setlength{\parindent}{0pt}
\setlength{\parskip}{5pt}
\setlength{\emergencystretch}{3em}
\newcommand{\N}{\mathbb{N}}
\newcommand{\Z}{\mathbb{Z}}
\newcommand{\Q}{\mathbb{Q}}
\newcommand{\R}{\mathbb{R}}
\newcommand{\C}{\mathbb{C}}
\newcommand{\F}{\mathbb{F}}
\newcommand{\A}{\mathbb{A}}
\newcommand{\PP}{\mathbb{P}}
\newcommand{\E}{\mathbb{E}}
\newcommand{\eps}{\varepsilon}
\newcommand{\dd}{\,\mathrm{d}}
\begin{document}
\section*{UnsolvedMath ID: 30004195}
\textbf{Source problem code:} catalogue entry
\subsection{Definitions and mathematical statement}
Let \(\R^3\) carry its Euclidean metric, and let \(\mathcal H^2\) denote two-dimensional Hausdorff measure. Define the constant-coefficient differential operator
\[
A=\partial_{x_1}^3+\partial_{x_2}^3+\partial_{x_3}^3.
\]
Interpret a real Radon measure \(\mu\) as a real signed Radon measure with locally finite total variation \(|\mu|\). It defines a distribution by integration against compactly supported smooth real-valued test functions. The condition \(A\mu=0\) in \(\mathcal D'(\R^3)\) means
\[
\int_{\R^3}\left(\partial_{x_1}^3\varphi+\partial_{x_2}^3\varphi+\partial_{x_3}^3\varphi\right)\,\dd\mu=0
\qquad\text{for every }\varphi\in C_c^\infty(\R^3;\R).
\]
Here \(\mu\ll\mathcal H^2\) means that every Borel set \(E\subseteq\R^3\) satisfying \(\mathcal H^2(E)=0\) also satisfies \(|\mu|(E)=0\).

\textbf{Question.} Is it true that every real Radon measure \(\mu\) on \(\R^3\) satisfies
\[
A\mu=0\text{ in }\mathcal D'(\R^3)
\quad\Longrightarrow\quad
\mu\ll\mathcal H^2?
\]
No global finiteness, support restriction, or boundary condition is specified in the supplied statement. The signed-measure interpretation of ``real Radon measure'' is made explicit here; the source record does not separately specify a positivity convention.
\subsection{Short English statement}
Must every real Radon measure on \(\R^3\) whose distributional third derivatives satisfy \(\partial_{x_1}^3\mu+\partial_{x_2}^3\mu+\partial_{x_3}^3\mu=0\) give zero total variation to every set of zero two-dimensional Hausdorff measure? The supplied record labels this question open; that status is not independently verified here.
\subsection{Sources}
\begin{itemize}
\item Local UnsolvedMath registry, supplied record with permanent ID \(30004195\), title ``Hausdorff Absolute Continuity of PDE-Constrained Measures,'' source problem code \texttt{catalogue entry}, in the Oberwolfach Reports Open Problems collection. The proposed year is recorded as 2019; no individual proposer is supplied.
\item Supplied source citation: \emph{Partial Differential Equations (2020)}, Oberwolfach Reports, DOI 10.4171/owr/2019/34, pp.\ 2033--2097. \url{https://doi.org/10.4171/owr/2019/34}. The supplied background describes these reports as documenting workshops and their open-problem discussions.
\item Supplied literature source, ``original Oberwolfach report'': \url{https://publications.mfo.de/bitstream/handle/mfo/3782/OWR_2019_34.pdf?isAllowed=y&sequence=1}.
\item Supplied literature source, \emph{Dimensional Estimates and Rectifiability for Measures Satisfying Linear PDE Constraints}: \url{https://doi.org/10.1007/s00039-019-00497-1}.
\item Literature assessment dated exactly \texttt{2026-08-22}: the supplied curation labels the problem open and reports that the cited dimension and rectifiability theory leaves this third-order operator conjecture unresolved, with no later proof or counterexample found. This is supplied dated curation derived from the OWR record, not an independently verified literature assessment. The record itself requests independent verification of time-sensitive status and cited claims.
\end{itemize}
\end{document}
